\documentclass[10pt,a4paper]{amsart}
\usepackage[margin=19mm]{geometry}
\usepackage{amsmath,amssymb,mathtools}
\usepackage[scr=rsfs,cal=euler]{mathalfa}
\usepackage{xfrac}
\usepackage[unicode,hidelinks,bookmarks=false]{hyperref}
\newcommand{\ie}{i.e.\ }
\newcommand{\cf}{cf.\ }
\newcommand{\resp}{respectively\ }
\newcommand{\Subsection}{\subsection}
\providecommand{\nobreakdash}{\nobreak\hbox{-}}
\setlength{\emergencystretch}{2em}
\multlinegap0pt
\usepackage{amssymb}
\usepackage{bbm}
\usepackage{xcolor}
\newtheorem{lemma}{Lemma}
\title{Strong maximum principle for fully nonlinear nonlocal problems}
\author{Juan Pablo Cabeza, Gabrielle Nornberg, Disson dos Prazeres}
\date{Source: arXiv:2602.13425v1 (CC BY 4.0)}
\begin{document}
\maketitle

\noindent\emph{Source and licence.} Strong maximum principle for fully nonlinear nonlocal problems. Version arXiv:2602.13425v1 (CC BY 4.0), distributed by arXiv under the Creative Commons Attribution 4.0 International licence, http://creativecommons.org/licenses/by/4.0/, as recorded in the arXiv licence metadata for this exact version and shown on the abstract page https://arxiv.org/abs/2602.13425v1. Full licence notice: \texttt{licenses/arxiv-2602.13425-CC-BY-4.0.txt}.\par\smallskip

\noindent\textit{Editorial scope.} Counted unit: the article's lemma supersolucion (source0.tex lines 204--215). The article's complete original proof of it is delivered with it (source0.tex lines 217--240, one \texttt{proof} environment of 24 lines). No mathematics is written, repaired or re-derived: the statement and the proof are the article's own lines, in the article's own notation, and no retained line is substituted or re-flowed.  Retained uncounted as the source-local context the proof needs: lines 188--200.  Nothing was omitted from any retained span, so the package carries no omission record.  No retained line contains a citation, so the wrapper carries no bibliography and no bibliography entry.  No retained line contains a figure environment, an image-inclusion command or drawing code, so no image is delivered and the package declares no asset dependency.  Editorial formatting only, in addition: an AMS \texttt{amsart} preamble that restates the article's own declarations and macros used by the retained lines, namely the environments \texttt{lemma} on separate counters, together with the standard packages those lines need.  The retained environments print in this document as Lemma 1 in order of appearance, whereas the article prints the same items with the numbering of its own full document.  The complete native source of the article as distributed in the arXiv e-print for this exact version is bundled unchanged as \texttt{sources/194636/source0.tex}.
\begin{lemma}\label{rho}
    Let \(s_0\in(0,1)\), \(s\in[s_0,1)\), and define
\begin{align*}
    \rho_{1}(x) = (\mathrm{dist}(x, \mathbb{R}^{n}\backslash B_{1}))^{s},  
    \quad
    \rho_{2}(x) = (\mathrm{dist}(x, \mathbb{R}^{n}\backslash B_{1}))^{3s/2}. 
\end{align*} Then, for all \(x\in B_{1}\backslash B_{1/2}\) we have
    \begin{align*}
        0  \geq \mathcal{M}^{-}[\rho_{1}](x) &\geq -C\{1+(1-s)|\log(1-|x|)|\}, \\
        \mathcal{M}^{-}[\rho_{2}](x) &\geq c(1-|x|)^{-s/2} - C,
    \end{align*}
    where the constants \(c>0\) and \(C\) depend only on \(n\), \(s_0\) and the ellipticity constants.
\end{lemma}

\begin{lemma}\label{supersolucion}
Let $s_0\in(0,1)$ and $s\in[s_0,1)$. There exists a radial  bounded continuous function $\varphi$ so that
\begin{equation}\label{rosoton1}
\begin{cases}
    \mathcal{M}^{-}[\varphi](x) - V^{-}(x)\varphi(x) \geq c & \text{ in } B_{1} \backslash B_{1/2}, \\ 
    \varphi(x) = 0 & \text{ in } \mathbb{R}^{n} \backslash B_{1}, \\
    \varphi(x) \geq c(1-|x|)^{s} & \text{ in } B_{1}, \\
    \varphi(x) \leq 1 & \text{ in } \overline{B_{1/2}},
\end{cases}
\end{equation}
for a positive constant $c$ that depends only on $n$, $s_0$ and the ellipticity constants.
\end{lemma}

\begin{proof}
    Consider the function \(\rho(x) :=\rho_{1}(x)+\rho_{2}(x)\); we proceed to construct a subsolution in the annulus \(B_{1}\backslash \overline{B_{1-\epsilon}}\). By Lemma~\ref{rho}, for \(|x|\in \left(\frac{1}{2},1\right)\) we have \[\mathcal{M}^{-}[\rho](x)-V^{-}(x)\rho(x)\geq -C\{1+(1-s)|\log(1-|x|)|\} + c(1-|x|)^{-s/2} - V^{-}(x)\left(\rho_{1}(x) + \rho_{2}(x)\right)-C,\]
    since \(\rho_{1}(x)=(1-|x|)^{s}\) and \(\rho_{2}(x) = (1-|x|)^{3s/2}\). Observing that both \(\rho_1\) and \(\rho_2\) vanish outside \(B_1\), and \(V^{-}\in L^{\infty}(B_{1})\), we conclude that \(-V^{-}(x)\rho(x)\geq - \|V^{-}\|_{\infty}\displaystyle\sup _{x\in B_{1}}\rho(x) = -2\|V^{-}\|_{\infty}\). It follows that
    \begin{equation*}
        \mathcal{M}^{-}[\rho](x)-V^{-}(x)\rho(x) \geq -C\{1+(1-s)|\log(1-|x|)|\} + c(1-|x|)^{-s/2} - 2\|V^{-}\|_{\infty} - C.
    \end{equation*}
    Hence, we can take \(\epsilon>0\) small enough such that \(\mathcal{M}^{-}[\rho] - V^{-}(x)\rho\geq1\) in \(B_{1}\backslash\overline{B_{1-\epsilon}}\), namely 
    \begin{equation*}
        \epsilon > \left( C_{1} + C_{2}|\log(1-|x|)|\right)^{-2/s},
    \end{equation*}
    with \(C_{1}:=c^{-1}\left(1+2C+2\|V^{-}\|_{\infty}\right)\) and \(C_{2}:=c^{-1}(1-s)C\).

    Given a large integer \(N\) and \(C>1\), we define \(0\leq k\leq N\), \(\displaystyle\widetilde{\rho}(x):=\max_{0\leq k\leq N} C^{k}\rho(2^{k/N}x)\) and \[A_{k} = \{x\in B_{1}:\widetilde{\rho}(x) = C^{k}\rho(2^{k/N}x)\}.\] 
    Since \(A_0 \subset B_{1}\backslash \overline{B_{1-\epsilon}}\), it follows that \(\mathcal{M}^{-}[\widetilde{\rho}]-V^{-}(x)\widetilde{\rho} \geq 1\) in \(A_0\). Further, we have \(A_k = 2^{-k/N}A_0\). For any \(x\in A_k\), it holds that \(2^{k/N}x\in A_0\subset B_{1}\backslash \overline{B_{1 - \epsilon}}\), thus
    \begin{align*}
        \mathcal{M}^{-}[\rho(2^{k/N})](x) &= 2^{2ks/N}\mathcal{M}^{-}[\rho](2^{k/N}x) 
        \geq \mathcal{M}^{-}[\rho](2^{k/N}x).
    \end{align*}
    For any \(x\in A_k\) we obtain the following estimate
    \begin{equation*}
        \mathcal{M}^{-}[\widetilde{\rho}](x) - V^{-}(x)\widetilde{\rho}(x) \geq C^{k}(\mathcal{M}^{-}[\widetilde{\rho}](2^{k/N}x) - V^{-}(2^{k/N}x)) \geq 1.
    \end{equation*}
    Finally, we set \(\varphi = c\widetilde{\rho}\) with \(c>0\) sufficiently small so that \(\varphi\leq 1\) in \(\overline{B_{1/2}}\).
\end{proof}

\end{document}
